\chapter{Singular Values}\label{linear:ch:svd}

\begin{center}\itshape
The singular value decomposition is the Swiss Army knife of linear algebra.
\upshape(a saying among numerical analysts)
\end{center}

In 1873 Eugenio Beltrami, a geometer in Bologna, published a short paper in a journal for university students. He wanted
to simplify a bilinear form $f(x,y)=x^{\mathsf T}Ay$ by changing $x$ and $y$ \emph{independently}, each by a rotation, and
showed that the form becomes $\sum_i\sigma_ix_iy_i$, the $\sigma_i^2$ being the eigenvalues of $A^{\mathsf T}A$. A year later
Camille Jordan, unaware of Beltrami, found the same thing in Paris. Sylvester rediscovered it in 1889, Erhard Schmidt in
1907 for integral operators, where he also proved that truncating the decomposition gives the best approximation of
lower rank. In 1936 the physicist Carl Eckart and the psychometrician Gale Young, writing in \emph{Psychometrika}, needed
exactly that theorem for the factor analysis of tables of test scores, proved it again for matrices, and their names
stuck to it.

The theorem then waited for computers. In 1965 Gene Golub and William Kahan showed how to compute singular values stably,
and the decomposition became the workhorse of statistics, signal processing and data analysis. Its most famous public
appearance was the Netflix Prize (2006--2009): a million dollars for predicting how users would rate films. The ratings
form a huge matrix with most entries missing, and the leading idea, popularised on a blog by a contestant writing as Simon
Funk, was to approximate it by a matrix of low rank, a product of a ``user'' matrix and a ``film'' matrix. Every user and
every film becomes a short vector of hidden tastes, and a rating is an inner product. The winning team blended hundreds of
models, but low rank was at the heart of most of them.

This chapter proves the singular value decomposition (SVD) for every rectangular matrix, reads it geometrically, and
derives its consequences: the polar decomposition, the natural norms of matrices, the Eckart--Young theorem, the
pseudo-inverse, principal component analysis and the condition number. The numerical side, how to compute all this
quickly and stably, is a subject of numerical linear algebra.

\begin{think}
A grey-scale photograph of $1000\times1000$ pixels is a matrix with a million entries. Suppose you may store only
$50{,}000$ numbers. What would you keep? Is there a sense in which one choice is \emph{best}, and could you prove it? The
answer, Theorem~\ref{linear:thm:eckartyoung}, is as clean as one could hope: keep the first $25$ singular triples.
\end{think}

\section{The decomposition}

Throughout, $\mathbb K=\R$ or $\C$, with the standard inner products, and $A\in M_{m,n}(\mathbb K)$. Unlike eigenvalues,
singular values need no square matrix.

\begin{theorem}[singular value decomposition]\label{linear:thm:svd}
Let $A\in M_{m,n}(\mathbb K)$ have rank $r$. There are orthonormal bases $v_1,\dots,v_n$ of $\mathbb K^n$ and
$u_1,\dots,u_m$ of $\mathbb K^m$, and numbers $\sigma_1\ge\sigma_2\ge\dots\ge\sigma_r>0$, such that
\[
Av_i=\sigma_iu_i\quad(i\le r),\qquad Av_i=0\quad(i>r).
\]
Equivalently $A=U\Sigma V^*=\sum_{i=1}^r\sigma_iu_iv_i^*$, with $U$, $V$ unitary and $\Sigma\in M_{m,n}$ zero except for
$\Sigma_{ii}=\sigma_i$. The \emph{singular values} $\sigma_i$ are uniquely determined: $\sigma_1^2,\dots,\sigma_r^2$ are the
non-zero eigenvalues of $A^*A$, and also of $AA^*$.
\end{theorem}
\begin{proof}
The matrix $A^*A$ is Hermitian and positive semidefinite ($x^*A^*Ax=\|Ax\|^2\ge0$), so by the spectral theorem
(Theorem~\ref{linear:thm:spectral}) it has an orthonormal eigenbasis $v_1,\dots,v_n$ with eigenvalues $s_1\ge\dots\ge s_n\ge0$.
Since $\ker A^*A=\ker A$ (Theorem~\ref{linear:thm:lsq}), exactly $r$ of them are positive; set $\sigma_i=\sqrt{s_i}$ for $i\le r$.
For $i>r$, $\|Av_i\|^2=s_i=0$. For $i\le r$, let $u_i=Av_i/\sigma_i$. These are orthonormal:
\[
\langle u_i,u_j\rangle=\frac{\langle A^*Av_i,v_j\rangle}{\sigma_i\sigma_j}=\frac{s_i\delta_{ij}}{\sigma_i\sigma_j}=\delta_{ij}.
\]
Extend them to an orthonormal basis of $\mathbb K^m$. Then $AV=U\Sigma$, which is the matrix form, and expanding $U\Sigma V^*$
column by row gives the sum of rank-one terms. Conversely, from any such decomposition, $A^*A=V\Sigma^*\Sigma V^*$ and
$AA^*=U\Sigma\Sigma^*U^*$, whose non-zero eigenvalues are the $\sigma_i^2$.
\end{proof}

The vectors $u_i$ and $v_i$ are the left and right \emph{singular vectors}. They are not unique (signs, and rotations inside
spaces with repeated singular values), but the numbers are. The decomposition gives orthonormal bases of the four
fundamental subspaces at once: $u_1,\dots,u_r$ span $\operatorname{Im}A$, $u_{r+1},\dots,u_m$ span $\ker A^*$, $v_1,\dots,v_r$
span $\operatorname{Im}A^*$ and $v_{r+1},\dots,v_n$ span $\ker A$, recovering Proposition~\ref{linear:prop:adjoint}.

\paragraph{Geometry.} Every linear map is a rotation, then a stretch along the axes, then another rotation. In particular,
the image of the unit sphere is an ellipsoid (flattened if $r<m$) whose semi-axes are the vectors $\sigma_iu_i$: for
$x=\sum c_iv_i$ with $\sum|c_i|^2=1$, $Ax=\sum_{i\le r}\sigma_ic_iu_i$, and the coordinates $y_i=\sigma_ic_i$ satisfy $\sum
y_i^2/\sigma_i^2\le1$.

\begin{figure}[ht]
\centering
\fig{linear-ch06-ellipse}
\caption{The matrix $A=\bigl(\begin{smallmatrix}3&0\\4&5\end{smallmatrix}\bigr)$ maps the unit circle to an ellipse with
semi-axes $\sigma_1=3\sqrt5$ and $\sigma_2=\sqrt5$ (drawn at reduced scale). The orthogonal directions $v_1,v_2$ are sent to
the orthogonal axes $\sigma_1u_1,\sigma_2u_2$; no other pair of orthogonal directions stays orthogonal.}
\label{linear:fig:ellipse}
\end{figure}

For a normal matrix $A=WDW^*$, an SVD is obtained by writing each eigenvalue as $\lambda_i=|\lambda_i|e^{i\theta_i}$, so the
singular values are the $|\lambda_i|$. For other matrices, eigenvalues and singular values can differ wildly: the Jordan
block $J_n$ has all eigenvalues $0$ and singular values $1,\dots,1,0$. They are related only by inequalities
(Exercise~\ref{linear:ex:normsvd}).

A useful device turns singular values into eigenvalues. The \emph{Hermitian dilation} of $A$ is
\[
\mathcal H(A)=\begin{pmatrix}0&A\\A^*&0\end{pmatrix}\in M_{m+n}(\mathbb K).
\]

\begin{proposition}\label{linear:prop:dilation}
With $p=\min(m,n)$ and $\sigma_i=0$ for $r<i\le p$, the eigenvalues of $\mathcal H(A)$ are $\pm\sigma_1,\dots,\pm\sigma_p$ and
$|m-n|$ zeros. In particular $\lambda_k(\mathcal H(A))=\sigma_k(A)$ for $k\le p$.
\end{proposition}
\begin{proof}
$\mathcal H(A)$ maps $(u,v)$ to $(Av,A^*u)$. For $i\le p$, the vectors $(u_i,\pm v_i)$ are eigenvectors with eigenvalues
$\pm\sigma_i$, since $Av_i=\sigma_iu_i$ and $A^*u_i=\sigma_iv_i$ (also when $\sigma_i=0$: then $u_i\perp\operatorname{Im}A$, so
$u_i\in\ker A^*$). The vectors $(u_j,0)$ for $j>p$ and $(0,v_j)$ for $j>p$ are in the kernel. These $2p+(m-p)+(n-p)=m+n$
vectors are orthogonal and non-zero, hence a basis.
\end{proof}

Through the dilation, everything proved about eigenvalues of Hermitian matrices in Chapter~\ref{linear:ch:spectral} transfers. For
example, Weyl's inequalities (Theorem~\ref{linear:thm:weyl}) for $\mathcal H(A)+\mathcal H(B)=\mathcal H(A+B)$ give, for $i+j-1\le
p$,
\begin{equation}\label{linear:eq:weylsv}
\sigma_{i+j-1}(A+B)\le\sigma_i(A)+\sigma_j(B).
\end{equation}

\section{The polar decomposition}

A complex number is $z=e^{i\theta}|z|$. Matrices factor the same way.

\begin{theorem}[polar decomposition]\label{linear:thm:polar}
Every $A\in M_n(\mathbb K)$ can be written $A=WP$ with $W$ unitary and $P$ positive semidefinite. The factor
$P=\sqrt{A^*A}$ is unique, and $W$ is unique when $A$ is invertible.
\end{theorem}
\begin{proof}
From $A=U\Sigma V^*$, write $A=(UV^*)(V\Sigma V^*)$: $UV^*$ is unitary and $V\Sigma V^*$ is positive semidefinite. If $A=WP$, then
$A^*A=PW^*WP=P^2$, so $P$ is the unique positive semidefinite square root of $A^*A$ (Exercise~\ref{linear:ex:sqrtpsd}). If $A$ is
invertible, so is $P$, and $W=AP^{-1}$.
\end{proof}

Geometrically, $P$ stretches along orthogonal axes and $W$ rotates (or reflects). Over $\R$ with $\det A>0$, the factor
$W$ is a rotation, and the path $W_tP_t$, with $P_t=(1-t)I+tP$ and $W_t$ a path of rotations from $I$ to $W$, connects $I$ to
$A$ inside the invertible matrices (a path of rotations exists by Exercise~\ref{linear:ex:realnormal}: turn each rotation block
gradually, and pair the $-1$ blocks, of which there is an even number, into rotations by $\pi$). So the real matrices of
positive determinant form a connected set.

\section{Norms of matrices}

The space $M_{m,n}(\mathbb K)$ has several natural norms, and singular values describe all of them.
\begin{itemize}
\item The \emph{operator norm} $\|A\|=\max_{\|x\|=1}\|Ax\|$ equals $\sigma_1$, since $\|Ax\|^2=\sum\sigma_i^2|c_i|^2$ for $x=
\sum c_iv_i$. It satisfies $\|AB\|\le\|A\|\|B\|$.
\item The \emph{Frobenius norm} $\|A\|_F=\bigl(\sum_{i,j}|a_{ij}|^2\bigr)^{1/2}=\sqrt{\operatorname{tr}A^*A}$ equals
$\bigl(\sum\sigma_i^2\bigr)^{1/2}$, since $\operatorname{tr}A^*A$ is the sum of the eigenvalues of $A^*A$.
\item The \emph{Ky Fan $k$-norm} is $\|A\|_{(k)}=\sigma_1+\dots+\sigma_k$; for $k=p$ it is the \emph{nuclear} (or trace) norm
$\|A\|_*=\sum\sigma_i$.
\end{itemize}
All of these are \emph{unitarily invariant}: $\|XAY\|=\|A\|$ for unitary $X,Y$, because $XAY$ has the same singular values.
That the Ky Fan functions are norms is not obvious, since singular values are not linear in $A$; here is the proof.

\begin{proposition}\label{linear:prop:kyfannorm}
For $k\le p$, $\|A+B\|_{(k)}\le\|A\|_{(k)}+\|B\|_{(k)}$.
\end{proposition}
\begin{proof}
By Proposition~\ref{linear:prop:dilation}, $\|A\|_{(k)}=\lambda_1(\mathcal H(A))+\dots+\lambda_k(\mathcal H(A))$. The sum of the $k$
largest eigenvalues is subadditive on Hermitian matrices (Exercise~\ref{linear:ex:kyfan}), and $\mathcal H$ is linear.
\end{proof}

\section{The best approximation of lower rank}

Truncate the SVD after $k$ terms: $A_k=\sum_{i\le k}\sigma_iu_iv_i^*$, a matrix of rank $k$ (for $k\le r$). Storing it costs
$k(m+n+1)$ numbers instead of $mn$.

\begin{theorem}[Schmidt, Eckart--Young, Mirsky]\label{linear:thm:eckartyoung}
For every matrix $B$ of rank at most $k$,
\[
\|A-B\|\ge\sigma_{k+1}=\|A-A_k\|\qquad\text{and}\qquad\|A-B\|_F^2\ge\sum_{i>k}\sigma_i^2=\|A-A_k\|_F^2.
\]
So $A_k$ is a best approximation of rank at most $k$ in both norms.
\end{theorem}
\begin{proof}
The equalities hold because $A-A_k=\sum_{i>k}\sigma_iu_iv_i^*$ is an SVD with singular values $\sigma_{k+1},\sigma_{k+2},\dots$.

Operator norm: $\ker B$ has dimension at least $n-k$, and $\operatorname{span}(v_1,\dots,v_{k+1})$ has dimension $k+1$, so they
share a unit vector $x=\sum_{i\le k+1}c_iv_i$. Then $\|(A-B)x\|=\|Ax\|=\bigl(\sum_{i\le k+1}\sigma_i^2|c_i|^2\bigr)^{1/2}\ge
\sigma_{k+1}$.

Frobenius norm: apply~\eqref{linear:eq:weylsv} to $A=(A-B)+B$ with $j=k+1$: since $\sigma_{k+1}(B)=0$,
$\sigma_{i+k}(A)\le\sigma_i(A-B)$ for every $i$ with $i+k\le p$. Squaring and summing over $i$,
$\sum_{i>k}\sigma_i(A)^2\le\sum_i\sigma_i(A-B)^2=\|A-B\|_F^2$.
\end{proof}

\begin{remark}
Mirsky (1960) proved that $A_k$ is a best approximation for \emph{every} unitarily invariant norm. The key is Ky Fan's
dominance theorem: $\|A\|\le\|B\|$ for all unitarily invariant norms if and only if $\|A\|_{(k)}\le\|B\|_{(k)}$ for every $k$
(see Bhatia's book in the sources); and~\eqref{linear:eq:weylsv} gives exactly these Ky Fan inequalities between $A-B$ and
$A-A_k$.
\end{remark}

\begin{figure}[ht]
\centering
\fig{linear-ch06-lowrank}
\caption{A $32\times32$ picture and its best approximations of rank $1$, $3$ and $8$. The relative errors in Frobenius norm,
$\bigl(\sum_{i>k}\sigma_i^2/\sum_i\sigma_i^2\bigr)^{1/2}$, are about $0.30$, $0.16$ and $0.07$. Rank $8$ needs $8\cdot65=520$
numbers instead of $1024$.}
\label{linear:fig:lowrank}
\end{figure}

For the photograph of the opening, $25$ singular triples cost $25\cdot2001\approx50{,}000$ numbers, and
Theorem~\ref{linear:thm:eckartyoung} says no other choice of $50{,}000$ numbers stored as a rank-$25$ matrix does better in
Frobenius or operator norm. (Whether the eye agrees is another question; JPEG uses other bases, adapted to the eye.)

A second reading of the theorem with $k=n-1$, for a square invertible $A$, is that \emph{$\sigma_n$ is the distance from $A$
to the nearest singular matrix}: the smallest singular value measures how close to non-invertible a matrix is.

\section{The pseudo-inverse and least squares}

When $A$ is not invertible, or not square, the SVD still inverts whatever can be inverted.

\begin{definition}\label{linear:def:pinv}
The \emph{Moore--Penrose pseudo-inverse} of $A=\sum_{i\le r}\sigma_iu_iv_i^*$ is $A^+=\sum_{i\le r}\sigma_i^{-1}v_iu_i^*\in
M_{n,m}(\mathbb K)$.
\end{definition}

\begin{theorem}\label{linear:thm:pinv}
For every $b\in\mathbb K^m$, $x^+=A^+b$ is the least squares solution of $Ax\approx b$ of smallest norm: it minimises $\|Ax-b\|$,
and among all minimisers it has the smallest $\|x\|$.
\end{theorem}
\begin{proof}
Write $x=\sum_jy_jv_j$ and $b=\sum_ic_iu_i$. Then
\[
\|Ax-b\|^2=\sum_{i\le r}|\sigma_iy_i-c_i|^2+\sum_{i>r}|c_i|^2.
\]
The minimisers are exactly the $x$ with $y_i=c_i/\sigma_i$ for $i\le r$ and $y_j$ arbitrary for $j>r$; among them,
$\|x\|^2=\sum|y_j|^2$ is smallest when $y_j=0$ for $j>r$, which is $x=\sum_{i\le r}\sigma_i^{-1}c_iv_i=A^+b$.
\end{proof}

When $A$ has independent columns, $A^+=(A^*A)^{-1}A^*$ and we recover Theorem~\ref{linear:thm:lsq}; when $A$ is invertible,
$A^+=A^{-1}$. Penrose (1955) characterised $A^+$ by four identities (Exercise~\ref{linear:ex:pinv}).

\section{Principal component analysis}

Let $x_1,\dots,x_N\in\R^n$ be data points, centred so that $\sum_ix_i=0$, forming the rows of $X\in M_{N,n}(\R)$. We look for
the $k$-dimensional subspace $W$ that fits them best, in the sense of the least total squared \emph{orthogonal} distance
$\sum_id(x_i,W)^2$. Pearson posed this in 1901, ``on lines and planes of closest fit''; Hotelling (1933) named the answer.

\begin{theorem}[principal components]\label{linear:thm:pca}
Let $X=\sum_i\sigma_iu_iv_i^{\mathsf T}$ be an SVD. The subspace $W_k=\operatorname{span}(v_1,\dots,v_k)$ minimises
$\sum_id(x_i,W)^2$ over all $k$-dimensional subspaces, with minimum $\sum_{j>k}\sigma_j^2$.
\end{theorem}
\begin{proof}
If $P_W$ is the orthogonal projection onto $W$, the rows of $XP_W$ are the projections of the data points, so
$\sum_id(x_i,W)^2=\|X-XP_W\|_F^2$, and $XP_W$ has rank at most $k$. By Theorem~\ref{linear:thm:eckartyoung} this is at least
$\sum_{j>k}\sigma_j^2$. For $W=W_k$, $XP_{W_k}=\sum_{i\le k}\sigma_iu_iv_i^{\mathsf T}=X_k$, which attains the bound.
\end{proof}

The $v_i$ are the \emph{principal directions}, the eigenvectors of the covariance matrix $\frac1NX^{\mathsf T}X$, and
$\sigma_i^2/N$ is the variance of the data along $v_i$. The fraction $\sum_{i\le k}\sigma_i^2/\sum_i\sigma_i^2$ is the
proportion of variance ``explained'' by $k$ components. Note the difference with least squares regression, which
minimises vertical distances and treats one coordinate as special: Exercise~\ref{linear:ex:pcatoy} shows the two lines differ, and
where the golden ratio hides.

\section{Conditioning}

How much does the solution of $Ax=b$ move when $b$ moves? For invertible $A$, the \emph{condition number} is
$\kappa(A)=\|A\|\|A^{-1}\|=\sigma_1/\sigma_n$.

\begin{proposition}\label{linear:prop:condition}
If $Ax=b\ne0$ and $A(x+\delta x)=b+\delta b$, then $\dfrac{\|\delta x\|}{\|x\|}\le\kappa(A)\dfrac{\|\delta b\|}{\|b\|}$, and equality
occurs for suitable $b$ and $\delta b$. Moreover $1/\kappa(A)$ is the relative distance from $A$ to the singular matrices:
$\min\{\|E\|/\|A\|:\ A+E\text{ singular}\}=\sigma_n/\sigma_1$.
\end{proposition}
\begin{proof}
$\delta x=A^{-1}\delta b$ gives $\|\delta x\|\le\sigma_n^{-1}\|\delta b\|$, and $\|b\|=\|Ax\|\le\sigma_1\|x\|$; multiply. Equality
holds for $x=v_1$ (so $b=\sigma_1u_1$) and $\delta b$ a multiple of $u_n$. The last statement is Theorem~\ref{linear:thm:eckartyoung}
with $k=n-1$, divided by $\|A\|=\sigma_1$.
\end{proof}

In floating point arithmetic with relative precision $10^{-16}$, a condition number of $10^{d}$ costs about $d$ digits, as
the Hilbert matrices of Chapter~\ref{linear:ch:inner} showed. Orthogonal matrices have $\kappa=1$, which is why algorithms built from
rotations and reflections are stable, and why $QR$ beats the normal equations, whose matrix $A^*A$ has condition number
$\kappa(A)^2$.

\section{Procrustes}

In Greek myth, Procrustes made his guests fit his bed by stretching or cutting them. The \emph{orthogonal Procrustes problem}
is gentler: given two configurations of points, the columns of $A,B\in M_{n,N}(\R)$ (say landmarks on two scans of a face),
find the rotation or reflection $Q$ that best superimposes them.

\begin{theorem}[orthogonal Procrustes]\label{linear:thm:procrustes}
Let $BA^{\mathsf T}=U\Sigma V^{\mathsf T}$ be an SVD. Then $Q=UV^{\mathsf T}$ minimises $\|QA-B\|_F$ over orthogonal $Q$.
\end{theorem}
\begin{proof}
$\|QA-B\|_F^2=\|A\|_F^2+\|B\|_F^2-2\operatorname{tr}(B^{\mathsf T}QA)$, and $\operatorname{tr}(B^{\mathsf T}QA)=\operatorname{tr}(QAB^{\mathsf T})
=\operatorname{tr}(QV\Sigma U^{\mathsf T})=\operatorname{tr}(Z\Sigma)$ with $Z=U^{\mathsf T}QV$ orthogonal. Since $|z_{ii}|\le1$
(the columns of $Z$ are unit vectors), $\operatorname{tr}(Z\Sigma)=\sum_i\sigma_iz_{ii}\le\sum_i\sigma_i$, with equality for
$Z=I$, that is $Q=UV^{\mathsf T}$.
\end{proof}

With $B=I$ and $A$ square, this says that the orthogonal matrix closest to $A$ is the factor $W$ of its polar decomposition
(Exercise~\ref{linear:ex:nearestorth}).

\begin{beauty}
Eigenvalues need a square matrix and a choice of basis that serves both sides of the map; singular values allow two
different orthonormal bases, one for the source and one for the target, and in exchange they always exist, are always real
and non-negative, and are perfectly stable. A single theorem then answers questions from statistics (which plane best fits a
cloud of points), numerical analysis (how close is a matrix to being singular), data compression (how few numbers describe a
picture), and geometry (what is the closest rotation): each is ``keep the largest singular values''.
\end{beauty}

\section{Exercises}

\begin{exercise}[$\star$]\label{linear:ex:svd2}
Compute an SVD of $A=\begin{pmatrix}3&0\\4&5\end{pmatrix}$ (Figure~\ref{linear:fig:ellipse}) and its polar decomposition. Check that
$\sigma_1\sigma_2=|\det A|$ and $\sigma_1^2+\sigma_2^2=\|A\|_F^2$.
\end{exercise}

\begin{exercise}[$\star$]\label{linear:ex:normsvd}
For square $A$, show that $|\det A|=\prod_i\sigma_i$ and $|\lambda|\le\sigma_1$ for every eigenvalue $\lambda$. Find a $2\times2$
matrix whose eigenvalues and singular values have nothing in common, and show that for a normal matrix the singular values
are the moduli of the eigenvalues.
\end{exercise}

\begin{exercise}[$\star$]\label{linear:ex:normineq}
Show that $\|A\|\le\|A\|_F\le\sqrt r\,\|A\|$, where $r=\operatorname{rank}A$, and that $\|AB\|_F\le\|A\|\,\|B\|_F$.
\end{exercise}

\begin{exercise}[$\star$]\label{linear:ex:pinv}
Show that $X=A^+$ satisfies Penrose's four equations $AXA=A$, $XAX=X$, $(AX)^*=AX$, $(XA)^*=XA$, that $AA^+$ and $A^+A$ are
the orthogonal projections onto $\operatorname{Im}A$ and $\operatorname{Im}A^*$, and compute the pseudo-inverse of
$\bigl(\begin{smallmatrix}1&1\\1&1\end{smallmatrix}\bigr)$.
\end{exercise}

\begin{exercise}[$\star\star$]\label{linear:ex:penroseunique}
Show that the four equations of Exercise~\ref{linear:ex:pinv} determine $X$ uniquely.
\end{exercise}

\begin{exercise}[$\star\star$]\label{linear:ex:pcatoy}
Find the first principal direction of the five points $(-2,-1)$, $(-1,-1)$, $(0,0)$, $(1,1)$, $(2,1)$, and compare with
the least squares line $y=bx$ through them. Where does the golden ratio $\varphi=\frac{1+\sqrt5}2$ appear?
\end{exercise}

\begin{exercise}[$\star\star$]\label{linear:ex:nearestorth}
Show that for square real $A$ with polar decomposition $A=WP$, the orthogonal matrix closest to $A$ in Frobenius norm is
$W$, and that $\|A-W\|_F^2=\sum_i(\sigma_i-1)^2$.
\end{exercise}

\begin{exercise}[$\star\star$]\label{linear:ex:nuclear}
Show that $\|A\|_*=\max|\operatorname{tr}(AQ)|$ over unitary $Q$ (for square $A$), and deduce $|\operatorname{tr}(AB)|\le\|A\|_*
\|B\|$.
\end{exercise}

\begin{exercise}[$\star\star$]\label{linear:ex:svinterlace}
Let $B$ be obtained from $A\in M_{m,n}$ by deleting one column. Show that $\sigma_1(A)\ge\sigma_1(B)\ge\sigma_2(A)\ge\sigma_2(B)
\ge\cdots$.
\end{exercise}

\begin{exercise}[$\star\star$]\label{linear:ex:svprod}
Show that $\sigma_i(AB)\le\|A\|\,\sigma_i(B)$ for all $i$. Deduce that $\kappa(AB)\le\kappa(A)\kappa(B)$ for invertible $A$,
$B$.
\end{exercise}

\begin{exercise}[$\star\star$]\label{linear:ex:normal2x2}
Show that the normal equations square the condition number: $\kappa(A^*A)=\kappa(A)^2$ for $A$ with independent columns
(with $\kappa=\sigma_1/\sigma_n$ for rectangular matrices), while $\kappa(R)=\kappa(A)$ for $A=QR$.
\end{exercise}

\begin{exercise}[$\star\star\star$]\label{linear:ex:matrixcompletion}
(A taste of Netflix.) A $3\times3$ matrix of rank $1$ has entries $a_{11}=2$, $a_{12}=4$, $a_{21}=3$, $a_{33}=5$, and
$a_{23}=10$; the others are unknown. Find all completions of rank $1$. Then show that a rank-one $m\times n$ matrix with all
entries non-zero is determined by the entries of one row, one column and nothing else, and that in general one needs
at least $m+n-1$ observed entries.
\end{exercise}

\section{Hints}

\begin{hint}{linear:ex:svd2}
$A^{\mathsf T}A=\bigl(\begin{smallmatrix}25&20\\20&25\end{smallmatrix}\bigr)$.
\end{hint}
\begin{hint}{linear:ex:normsvd}
$|\det A|^2=\det A^*A$. For the eigenvalue, take a unit eigenvector. Try a Jordan block with a non-zero eigenvalue.
\end{hint}
\begin{hint}{linear:ex:normineq}
Use $\|A\|_F^2=\sum\sigma_i^2$; for the product, apply $A$ to each column of $B$.
\end{hint}
\begin{hint}{linear:ex:pinv}
Check everything in the singular bases, where $A$ and $A^+$ are ``diagonal''. The matrix $\bigl(\begin{smallmatrix}1&1\\1&1\end{smallmatrix}\bigr)$
is $2uu^{\mathsf T}$ with $u=(1,1)/\sqrt2$.
\end{hint}
\begin{hint}{linear:ex:penroseunique}
If $X$ and $Y$ both satisfy the equations, write $X=XAX=X(AX)^*=\dots$ and move factors around until $Y$ appears.
\end{hint}
\begin{hint}{linear:ex:pcatoy}
$X^{\mathsf T}X=\bigl(\begin{smallmatrix}10&6\\6&4\end{smallmatrix}\bigr)$.
\end{hint}
\begin{hint}{linear:ex:nearestorth}
Theorem~\ref{linear:thm:procrustes} with $B=I$; then compute in the singular bases.
\end{hint}
\begin{hint}{linear:ex:nuclear}
$\operatorname{tr}(AQ)=\operatorname{tr}(\Sigma V^*QU)$, as in the proof of Theorem~\ref{linear:thm:procrustes}.
\end{hint}
\begin{hint}{linear:ex:svinterlace}
$B^*B$ is a principal submatrix of $A^*A$.
\end{hint}
\begin{hint}{linear:ex:svprod}
$\|ABx\|\le\|A\|\|Bx\|$, and Courant--Fischer for $B^*A^*AB$ and $B^*B$.
\end{hint}
\begin{hint}{linear:ex:normal2x2}
$A^*A=V\Sigma^*\Sigma V^*$, and $A^*A=R^*R$.
\end{hint}
\begin{hint}{linear:ex:matrixcompletion}
A rank-one matrix is $xy^{\mathsf T}$, so $a_{ij}a_{kl}=a_{il}a_{kj}$.
\end{hint}

\section{Solutions}

\begin{solution}{linear:ex:svd2}
\idea Diagonalise $A^{\mathsf T}A$ to get $V$ and the $\sigma_i$; then $u_i=Av_i/\sigma_i$.
\writeup $A^{\mathsf T}A$ has eigenvalues $45$ and $5$ with eigenvectors $v_1=\frac1{\sqrt2}(1,1)$, $v_2=\frac1{\sqrt2}(-1,1)$.
So $\sigma_1=3\sqrt5$, $\sigma_2=\sqrt5$, and $u_1=\frac{Av_1}{\sigma_1}=\frac{(3,9)}{\sqrt2\cdot3\sqrt5}=\frac1{\sqrt{10}}(1,3)$,
$u_2=\frac{Av_2}{\sigma_2}=\frac{(-3,1)}{\sqrt2\sqrt5}=\frac1{\sqrt{10}}(-3,1)$. Then $\sigma_1\sigma_2=15=|\det A|$ and
$\sigma_1^2+\sigma_2^2=50=9+16+25$. For the polar decomposition, $P=\sigma_1v_1v_1^{\mathsf T}+\sigma_2v_2v_2^{\mathsf T}
=\frac{3\sqrt5}2\bigl(\begin{smallmatrix}1&1\\1&1\end{smallmatrix}\bigr)+\frac{\sqrt5}2\bigl(\begin{smallmatrix}1&-1\\-1&1\end{smallmatrix}\bigr)
=\sqrt5\bigl(\begin{smallmatrix}2&1\\1&2\end{smallmatrix}\bigr)$, and
$W=UV^{\mathsf T}=u_1v_1^{\mathsf T}+u_2v_2^{\mathsf T}=\frac1{\sqrt5}\bigl(\begin{smallmatrix}2&-1\\1&2\end{smallmatrix}\bigr)$, a
rotation by $\arctan\frac12$. Check: $WP=\bigl(\begin{smallmatrix}2&-1\\1&2\end{smallmatrix}\bigr)\bigl(\begin{smallmatrix}2&1\\1&2
\end{smallmatrix}\bigr)=\bigl(\begin{smallmatrix}3&0\\4&5\end{smallmatrix}\bigr)$.
\end{solution}

\begin{solution}{linear:ex:normsvd}
\idea Determinants and norms are unitarily invariant, so compute them on $\Sigma$.
\writeup $|\det A|^2=\det(A^*A)=\prod\sigma_i^2$. If $Ax=\lambda x$ with $\|x\|=1$, then $|\lambda|=\|Ax\|\le\sigma_1$. The matrix
$\bigl(\begin{smallmatrix}1&1\\0&1\end{smallmatrix}\bigr)$ has eigenvalues $1,1$ and $A^{\mathsf T}A=\bigl(\begin{smallmatrix}1&1\\1&2
\end{smallmatrix}\bigr)$ with eigenvalues $\frac{3\pm\sqrt5}2$, so $\sigma_{1,2}=\frac{\sqrt5\pm1}2=\varphi^{\pm1}$: neither
is $1$. For normal $A=WDW^*$, write $D=|D|\Theta$ with $\Theta$ diagonal unitary; then $A=(W\Theta)|D|W^*$ is an SVD (after
reordering), with singular values $|\lambda_i|$.
\end{solution}

\begin{solution}{linear:ex:normineq}
\idea Compare $\max\sigma_i$, $\bigl(\sum\sigma_i^2\bigr)^{1/2}$ over $r$ non-zero terms.
\writeup $\sigma_1^2\le\sum_{i\le r}\sigma_i^2\le r\sigma_1^2$. For the product, if $b_j$ are the columns of $B$, then
$\|AB\|_F^2=\sum_j\|Ab_j\|^2\le\|A\|^2\sum_j\|b_j\|^2=\|A\|^2\|B\|_F^2$.
\end{solution}

\begin{solution}{linear:ex:pinv}
\idea In the bases $(v_j)$ and $(u_i)$, $A$ is $\Sigma$ and $A^+$ is $\Sigma^+$, the transpose of $\Sigma$ with non-zero entries
inverted; everything reduces to diagonal matrices.
\writeup $AA^+=\sum_{i\le r}u_iu_i^*$ and $A^+A=\sum_{i\le r}v_iv_i^*$, the orthogonal projections onto $\operatorname{span}
(u_1,\dots,u_r)=\operatorname{Im}A$ and $\operatorname{span}(v_1,\dots,v_r)=\operatorname{Im}A^*$. They are Hermitian, and
$AA^+A=\bigl(\sum_{i\le r}u_iu_i^*\bigr)A=A$, $A^+AA^+=\bigl(\sum_{i\le r}v_iv_i^*\bigr)A^+=A^+$. For $A=2uu^{\mathsf T}$ (an SVD
with $\sigma_1=2$, $u_1=v_1=u$), $A^+=\frac12uu^{\mathsf T}=\frac14\bigl(\begin{smallmatrix}1&1\\1&1\end{smallmatrix}\bigr)$.
\end{solution}

\begin{solution}{linear:ex:penroseunique}
\idea Use the equations to insert $A$ between factors, and the two Hermitian conditions to move projections across.
\writeup Let $X,Y$ satisfy the four equations. Then $AX=(AX)^*=X^*A^*=X^*(AYA)^*=X^*A^*Y^*A^*=(AX)^*(AY)^*=AXAY=AY$, and
symmetrically $XA=(XA)^*=A^*X^*=(AYA)^*X^*=A^*Y^*A^*X^*=(YA)^*(XA)^*=YAXA=YA$. Hence $X=XAX=X(AY)=(XA)Y=(YA)Y=Y$.
\end{solution}

\begin{solution}{linear:ex:pcatoy}
\idea PCA diagonalises $X^{\mathsf T}X$; regression solves one normal equation. Both are two-line computations, and the
golden ratio comes from $\sqrt{45}=3\sqrt5$.
\writeup The points are centred. $X^{\mathsf T}X=\bigl(\begin{smallmatrix}10&6\\6&4\end{smallmatrix}\bigr)$ has trace $14$ and
determinant $4$, so eigenvalues $7\pm3\sqrt5$. An eigenvector for $7+3\sqrt5$ is $(6,\,3\sqrt5-3)$, of slope $\frac{\sqrt5-1}2=
\frac1\varphi\approx0.618$. The least squares line $y=bx$ has $b=\frac{\sum xy}{\sum x^2}=\frac6{10}=0.6$, a little flatter:
regression minimises vertical distances and so is pulled towards the $x$-axis, while PCA treats both coordinates alike. The
fraction of variance explained by the first component is $\frac{7+3\sqrt5}{14}\approx0.979$.
\end{solution}

\begin{solution}{linear:ex:nearestorth}
\idea Procrustes with target $I$: minimise $\|QA-I\|_F=\|A-Q^{\mathsf T}\|_F$.
\writeup By Theorem~\ref{linear:thm:procrustes} with $B=I$, and $BA^{\mathsf T}=A^{\mathsf T}=V\Sigma U^{\mathsf T}$, the minimiser of
$\|QA-I\|_F$ is $Q=VU^{\mathsf T}$. Since $\|QA-I\|_F=\|A-Q^{\mathsf T}\|_F$ for orthogonal $Q$, the closest orthogonal matrix to
$A$ is $Q^{\mathsf T}=UV^{\mathsf T}=W$. Finally $A-W=U(\Sigma-I)V^{\mathsf T}$ has Frobenius norm $\bigl(\sum(\sigma_i-1)^2\bigr)^{1/2}$.
\end{solution}

\begin{solution}{linear:ex:nuclear}
\idea The trace of $\Sigma$ times a unitary matrix is a combination of the $\sigma_i$ with coefficients of modulus at most $1$.
\writeup $\operatorname{tr}(AQ)=\operatorname{tr}(U\Sigma V^*Q)=\operatorname{tr}(\Sigma Z)$ with $Z=V^*QU$ unitary, so
$|\operatorname{tr}(AQ)|=|\sum\sigma_iz_{ii}|\le\sum\sigma_i$, with equality for $Q=VU^*$. For general $B$ with $\|B\|\le1$:
$B=W P$ with $P=\sum_j\tau_jw_jw_j^*$, $0\le\tau_j\le1$, is a convex combination of unitary matrices (write each $\tau_j=
\frac{e^{i\theta_j}+e^{-i\theta_j}}2$ with $\cos\theta_j=\tau_j$, so $P=\frac{U_++U_-}2$ with $U_\pm=\sum_je^{\pm i\theta_j}w_jw_j^*$
unitary). Hence $|\operatorname{tr}(AB)|\le\frac12(|\operatorname{tr}(AWU_+)|+|\operatorname{tr}(AWU_-)|)\le\|A\|_*$, and scaling
gives $|\operatorname{tr}(AB)|\le\|A\|_*\|B\|$.
\end{solution}

\begin{solution}{linear:ex:svinterlace}
\idea Squares of singular values are eigenvalues of Gram matrices, and deleting a column deletes a row and a column of
the Gram matrix.
\writeup If $B$ is $A$ without column $j$, then $B^*B$ is $A^*A$ without row and column $j$. By Cauchy's interlacing
(Theorem~\ref{linear:thm:interlacing}), $\lambda_k(A^*A)\ge\lambda_k(B^*B)\ge\lambda_{k+1}(A^*A)$, and taking square roots (all are
$\ge0$) gives the claim.
\end{solution}

\begin{solution}{linear:ex:svprod}
\idea $\|ABx\|\le\|A\|\|Bx\|$ compares the Rayleigh quotients of $(AB)^*(AB)$ and $\|A\|^2B^*B$.
\writeup For every $x$, $x^*(AB)^*ABx=\|ABx\|^2\le\|A\|^2\|Bx\|^2=\|A\|^2x^*B^*Bx$. By Courant--Fischer
(Theorem~\ref{linear:thm:courantfischer}), $\lambda_i\bigl((AB)^*AB\bigr)\le\|A\|^2\lambda_i(B^*B)$, that is $\sigma_i(AB)\le\|A\|\sigma_i(B)$.
By symmetry ($\sigma_i(AB)=\sigma_i(B^*A^*)$) also $\sigma_i(AB)\le\sigma_i(A)\|B\|$. For invertible $A,B$ of size $n$:
$\sigma_1(AB)\le\sigma_1(A)\sigma_1(B)$, and $\sigma_n(AB)=1/\|(AB)^{-1}\|\ge1/(\|B^{-1}\|\|A^{-1}\|)=\sigma_n(A)\sigma_n(B)$.
\end{solution}

\begin{solution}{linear:ex:normal2x2}
\idea Singular values of $A^*A$ are the squares of those of $A$; those of $R$ equal those of $A$ because $Q$ is an isometry.
\writeup $A^*A=V(\Sigma^*\Sigma)V^*$ is an SVD (and spectral decomposition) with singular values $\sigma_i^2$, so $\kappa(A^*A)=
\sigma_1^2/\sigma_n^2$. For $A=QR$ with $Q^*Q=I$, $A^*A=R^*R$, so $R$ and $A$ have the same singular values and $\kappa(R)=
\kappa(A)$. Solving $Rx=Q^*b$ therefore loses about half as many digits as solving the normal equations.
\end{solution}

\begin{solution}{linear:ex:matrixcompletion}
\idea Rank one means all $2\times2$ minors vanish: $a_{ij}a_{kl}=a_{il}a_{kj}$. Each known pair of entries in a common row or
column propagates information.
\writeup Write $A=xy^{\mathsf T}$. From $a_{11}=x_1y_1=2$ and $a_{12}=x_1y_2=4$, $y_2=2y_1$; from $a_{21}=x_2y_1=3$,
$x_2=\frac32x_1$. From $a_{23}=x_2y_3=10$, $y_3=\frac{20}{3x_1}$, so $a_{13}=x_1y_3=\frac{20}3$ and $a_{22}=x_2y_2=6$. From
$a_{33}=x_3y_3=5$, $x_3=\frac{3x_1}4$, so $a_{31}=x_3y_1=\frac34\cdot2=\frac32$ and $a_{32}=3$. All unknowns are determined:
the completion is unique,
\[
A=\begin{pmatrix}2&4&\frac{20}3\\3&6&10\\\frac32&3&5\end{pmatrix}=\begin{pmatrix}1\\\frac32\\\frac34\end{pmatrix}\begin{pmatrix}2&4&\frac{20}3\end{pmatrix}.
\]
In general, with all entries non-zero, row $i$ and column $j$ determine $a_{kl}=a_{kj}a_{il}/a_{ij}$. Counting: the set of
rank-one matrices is described by $x\in\mathbb K^m$, $y\in\mathbb K^n$ up to the scaling $(tx,t^{-1}y)$, that is $m+n-1$
parameters; more precisely, the bipartite graph of observed positions (rows and columns as vertices) must be connected for
uniqueness, since otherwise one connected component can be rescaled independently, and a connected graph on $m+n$ vertices
has at least $m+n-1$ edges. Matrix completion of rank $k$ from random entries, the mathematics behind the Netflix Prize, is
the subject of Cand\`es and Recht's paper in the sources.
\end{solution}

\section*{Sources and further reading}

\begin{itemize}
\item E.~Beltrami, \emph{Sulle funzioni bilineari}, Giornale di Matematiche ad uso degli studenti delle universit\`a 11
(1873), 98--106; C.~Jordan, \emph{M\'emoire sur les formes bilin\'eaires}, Journal de math\'ematiques pures et appliqu\'ees
19 (1874), 35--54.
\item G.~W.~Stewart, \emph{On the early history of the singular value decomposition}, SIAM Review 35 (1993), 551--566.
\item E.~Schmidt, \emph{Zur Theorie der linearen und nichtlinearen Integralgleichungen I}, Mathematische Annalen 63 (1907),
433--476; C.~Eckart and G.~Young, \emph{The approximation of one matrix by another of lower rank}, Psychometrika 1 (1936),
211--218; L.~Mirsky, \emph{Symmetric gauge functions and unitarily invariant norms}, Quarterly Journal of Mathematics 11
(1960), 50--59.
\item K.~Pearson, \emph{On lines and planes of closest fit to systems of points in space}, Philosophical Magazine 2 (1901),
559--572; H.~Hotelling, \emph{Analysis of a complex of statistical variables into principal components}, Journal of
Educational Psychology 24 (1933), 417--441.
\item R.~Penrose, \emph{A generalized inverse for matrices}, Proceedings of the Cambridge Philosophical Society 51 (1955),
406--413.
\item G.~Golub and W.~Kahan, \emph{Calculating the singular values and pseudo-inverse of a matrix}, SIAM Journal on
Numerical Analysis 2 (1965), 205--224.
\item P.~H.~Sch\"onemann, \emph{A generalized solution of the orthogonal Procrustes problem}, Psychometrika 31 (1966), 1--10.
\item Y.~Koren, R.~Bell and C.~Volinsky, \emph{Matrix factorization techniques for recommender systems}, IEEE Computer 42
(2009), 30--37; E.~J.~Cand\`es and B.~Recht, \emph{Exact matrix completion via convex optimization}, Foundations of
Computational Mathematics 9 (2009), 717--772.
\item R.~Bhatia, \emph{Matrix Analysis}, Springer, 1997, Chapter~IV (unitarily invariant norms, Ky Fan dominance); L.~N.~Trefethen
and D.~Bau, \emph{Numerical Linear Algebra}, SIAM, 1997, Lectures 4--5.
\end{itemize}
